% ============================================================================
%  SECCIÓN 3.6: FLUJO INICIAL DE NAVEGACIÓN
% ============================================================================

\section{Flujo inicial de navegaci\'on}

El flujo de navegaci\'on describe el recorrido que realiza el usuario a trav\'es
de las pantallas de la aplicaci\'on para cumplir sus tareas principales
\citep{pressman2010}. A continuaci\'on se describe el recorrido principal y el
diagrama de navegaci\'on propuesto.

\subsection{Recorrido principal}

Al abrir la aplicaci\'on, el usuario visualiza la pantalla de bienvenida
(splash) y puede registrarse o iniciar sesi\'on. Una vez autenticado, accede a
la pantalla principal, que combina el mapa interactivo con el cat\'alogo de
escenarios deportivos. Desde esta pantalla puede seleccionar un escenario y
consultar su detalle, as\'i como guardarlo en sus favoritos para acceder a \'el
de forma r\'apida en visitas posteriores.

\subsection{Diagrama de navegaci\'on}

\begin{figure}[H]
  \centering
  \caption{Flujo de navegaci\'on propuesto para la aplicaci\'on m\'ovil}
  \contenidoConFuente{%
  \begin{tikzpicture}[
      nodo/.style={rectangle, draw=black!70, rounded corners=2pt, minimum width=3.4cm,
                   minimum height=0.95cm, align=center, fill=gray!10, font=\small},
      flecha/.style={-{Stealth[length=2.5mm]}, thick}
    ]
    \node[nodo] (splash) at (0,0) {Splash / Bienvenida};
    \node[nodo] (registro) at (-3.4,-2.0) {Registro de usuario};
    \node[nodo] (login) at (3.4,-2.0) {Inicio de sesi\'on};
    \node[nodo] (principal) at (0,-4.0) {Pantalla principal\\ (Mapa + Cat\'alogo)};
    \node[nodo] (detalle) at (-3.4,-6.0) {Detalle del escenario};
    \node[nodo] (favoritos) at (3.4,-6.0) {Favoritos};

    \draw[flecha] (splash) -- (registro);
    \draw[flecha] (splash) -- (login);
    \draw[flecha] (registro) -- (login);
    \draw[flecha] (login) -- (principal);
    \draw[flecha] (registro) -- (principal);
    \draw[flecha] (principal) -- (detalle);
    \draw[flecha] (principal) -- (favoritos);
    \draw[flecha] (detalle) -- (favoritos);
    \draw[flecha] (favoritos) -- (detalle);
  \end{tikzpicture}%
  }{Elaboraci\'on propia.}
\end{figure}
